\documentclass[11pt,a4paper]{amsart}
\usepackage[a4paper,margin=20mm]{geometry}
\usepackage{amsmath,amssymb,amsthm,graphicx,xcolor}
\usepackage[no-math]{fontspec}
\setmainfont{lmroman10-regular.otf}[BoldFont=lmroman10-bold.otf,ItalicFont=lmroman10-italic.otf,BoldItalicFont=lmroman10-bolditalic.otf]
\usepackage{longtable,booktabs,mathtools,mathrsfs,url,hyperref,xifthen,paralist,enumitem,bussproofs,subcaption,doi}
\usepackage[capitalise]{cleveref}
\AtBeginDocument{\def\shorttitle{Dependently sorted nominal signatures}}
\newtheorem{definition}{Definition}
\newtheorem{lemma}{Lemma}
\newtheorem{theorem}{Theorem}
\newcounter{ruleName}
\newcommand{\ruleNameLbl}[1]{\refstepcounter{ruleName}(#1)\label{rule:#1}}
\newcounter{s}
\setcounter{s}{-1}


\newcommand{\prg}[1]{\underline{#1}}

\newcommand{\atomsort}[1]{\ensuremath{\mathsf{#1}}}
\newcommand{\atomset}[2]{\ensuremath{{\mathbb{#1}}_{#2}}}
\newcommand{\atom}[1]{\ensuremath{#1}}
\newcommand{\swap}[2]{\ensuremath{(\atom{#1}\,\atom{#2})}}
\newcommand{\swapapp}[2]{#1 \cdot #2}
\newcommand{\datasort}[1]{\ensuremath{\mathrm{#1}}}
\newcommand{\absort}[2]{\ensuremath{\ll\!\!#1\!\!\gg #2}}
\newcommand{\absterm}[2]{\ensuremath{\ll\!\!#1\!\!\gg #2}}
\newcommand{\dabsort}[3]{\ensuremath{\ll\!\!\ntyping{#1}{#2}\!\!\gg #3}}
\newcommand{\dabsterm}[3]{\ensuremath{\ll\!\!\ntyping{#1}{#2}\!\!\gg #3}}
\newcommand{\concr}[2]{#1[#2]}
\newcommand{\concrVect}[2]{#1\vect{[#2]}}
\newcommand{\mvarsubs}[2]{#1 #2}
\newcommand{\vect}[1]{\ensuremath{\overline{#1}}}
\newcommand{\simpleconstrdecl}[3]{\ensuremath{#1 : #2 \rightarrow\ #3}}
\newcommand{\sconstrdecl}[3]{\ensuremath{#1 : (\_:#2) \rightarrow\ #3}}
\newcommand{\depconstrdecl}[3]{\ensuremath{#1 : (#2) \rightarrow\ #3}}
\newcommand{\sconstr}[2]{\ensuremath{#1\,#2}}
\newcommand{\typing}[2]{\ensuremath{#1 : #2}}
\newcommand{\ntyping}[2]{\ensuremath{#1{:}#2}}
\newcommand{\fresh}[2]{#1 \# #2}
\newcommand{\alphacnv}[3]{\ensuremath{#1 =_{\alpha} #2 : #3}}
\newcommand{\sort}{\ensuremath{\mathsf{sort}}}
\newcommand{\data}{\ensuremath{\mathsf{data}}}
\newcommand{\datsort}{\ensuremath{\mathsf{datasort}}}
\newcommand{\constrdecl}[3]{\ensuremath{(\vect{\typing{#1}{#2}}) \rightarrow\ #3}}
\newcommand{\dconstrdecl}[4]{\ensuremath{#1\,{:}\,(\vect{\typing{#2}{#3}}) \rightarrow\ #4}}
\newcommand{\sortconstrdecl}[2]{\ensuremath{#1 : #2 \rightarrow\ \data}}
\newcommand{\fconstrdecl}[4]{\ensuremath{#1 : #2 \rightarrow\ #3\ ;\ #4}}
\newcommand{\fconstrdef}[2]{\ensuremath{#1 \equiv #2}}

%% Commands for judgements.

\newcommand{\typeAndFresh}[3]{#1\colon #2\,;\,#3}
\newlength{\jdgsep}
\setlength{\jdgsep}{0.3em}

\newcommand{\emptyTel}{\cdot}
\newcommand{\consTel}[3]{#1,(\typing{#2}{#3})}
\newcommand{\emptyCtx}{\cdot}
\newcommand{\consCtx}[3]{#1,(\typing{#2}{#3})}
\newcommand{\emptySubst}{\cdot}
\newcommand{\consSubst}[3]{#1,[#2\mapsto #3]}

\newcommand{\lookupConstSig}[3]{#2(#1) = #3}

\newcommand{\sigjdgmnt}[1]{\ensuremath{\vdash #1\hspace{\jdgsep} \text{sig-ok}}}
\newcommand{\teljdgmntsig}[3]{\ensuremath{#1\vdash_{#2} #3 \hspace{\jdgsep} \text{tel-ok}}}
\newcommand{\telfitsjdgmntsig}[5]{\ensuremath{#1;#2\vdash_{#3} #4 \text{ fits } #5}}
\newcommand{\sortjdgmntsig}[5]{\ensuremath{#1;#2\vdash_{#4}  #5 \ \ \sort}}
\newcommand{\sortjdgmnt}[4]{\ensuremath{#1;#2;#3\vdash #4 \sort}}
\newcommand{\clsortjdgmntsig}[4]{\ensuremath{#1;\emptyCtx\vdash_{#3} #4 \ \ \sort}}
\newcommand{\ctxjdgmntsig}[4]{\ensuremath{#1\vdash_{#3} #4 \hspace{\jdgsep} \text{ctx-ok}}}
\newcommand{\varSubjdgmntsig}[6]{\ensuremath{#1;#2\vdash_{#3} \typing{#4}{#5 \to #6}}}
\newcommand{\atmSubjdgmntsig}[6]{\ensuremath{#1;#2\vdash_{#3} \typing{#4}{#5 \to #6}}}
\newcommand{\tyjdgmntsig}[6]{\ensuremath{#1;#2\vdash_{#4} \typing{#5}{#6}}}
\newcommand{\tyjdgmnt}[5]{\ensuremath{#1;#2\vdash \typing{#4}{#5}}}
\newcommand{\tyjdgmnts}[2]{\ensuremath{\vdash \typing{#1}{#2}}}
\newcommand{\sigtjdgmnt}[2]{\ensuremath{#1\vdash #2}}
\newcommand{\groundtyjdgmntsig}[4]{\ensuremath{#1\vdash_{#2} \typing{#3}{#4}}}

\newcommand{\SigEmpty}{\ensuremath{\langle \rangle}}
\newcommand{\SigConsSort}[4]{\ensuremath{#1 , \langle %
    \typeAndFresh{#2}{#3\to\data}{#4} \rangle}}
\newcommand{\SigConsFun}[5]{\ensuremath{#1 , \langle %
    \typeAndFresh{#2}{#3\to #4}{#5} \rangle}}

%%
%%
\newcommand{\ifthen}[4]{\ensuremath{\mathsf{if}\, #1 \, \mathsf{fresh\ in}\, #2\, \mathsf{then}\, #3\, \mathsf{else}\, #4}}
\newcommand{\ifthennn}[4]{\ensuremath{\mathsf{case}\, (\fresh{#1}{#2},#3,#4)}}
%% Comment the next line for the if_fresh in_then_else_ version 
\renewcommand{\ifthen}{\ifthennn}

\newcommand{\iif}[4]{\ensuremath{\mathsf{caseEval}\, (\fresh{#1}{#2},  #3, #4)}}


\newcommand{\cnst}[1]{\textsl{#1}}
\newcommand{\set}{\textsl{Set}}
\newcommand{\el}{\textsl{El}}
\newcommand{\form}{\texttt{Form}}
\newcommand{\prf}{\ensuremath{\mathcal{D}}}
\newcommand{\impl}{\ensuremath{\supset}}
\newcommand{\nat}{\ensuremath{\mathbb{N}}}
\newcommand{\eqs}{\ensuremath{\mathit{Eq}}}
\newcommand{\name}{\textsl{V}}
\newcommand{\ofType}{\colon}
\newcommand {\deffont}  [1] {\emph{#1}}
\newcommand {\id}       [0] {\mathsf{id}}
\newcommand {\act}      [0] {\cdot}
\newcommand {\prmInv}   [1] {#1^{-1}}
\newcommand {\prmAct}   [2] {#1 \act #2}
\newcommand {\prmAtm}    [2] {#1 \act #2}    %%FIXME: VERIFICAR ESTE COMANDO
\newcommand{\aSub}[3]{\ensuremath{[#1 \mapsto #2]#3}}
%\newcommand{\aSubAtm}[2]{\ensuremath{#1 #2}} %%FIXME: VERIFICAR ESTE COMANDO
\newcommand {\val}      [2] {#1({#2})}
\newcommand {\aSubst}   [2] {\map{#1}{#2}}
\newcommand {\aSubAtm} [2] {\val{#2}{#1}}
\newcommand {\aSubAct} [2] {#1 \, #2}
\newcommand{\atomsubs}[3]{\ensuremath{\aSubAct{#3}{[#1 \atomsubstsym #2]}}}
\newcommand{\atomsubstsym}{\ensuremath{\mapsto}}

\newcommand{\varTrmD} [3] {\ensuremath{#3}} %\susp{\ensuremath{#1}}{\ensuremath{#2}} \act \ensuremath{#3}}
\newcommand{\varConcr} [2] {{#1}[#2]}
\newcommand {\tplTrm}  [1] {(#1)}
\newcommand {\absTrmD}  [3] {\absterm{#1 \ofType #2} #3}
\newcommand {\appTrm}  [2] {\sconstr{#1}{#2}}
\newcommand {\absTypD}  [3] {\absort{#1 \ofType #2}{#3}}
\newcommand {\prdTypD}  [1] {(#1)}
\newcommand {\cnsTyp}  [2] {\sconstr{#1}{#2}}
\newcommand {\susp}    [2] {\ifthenelse{\isempty{#1}}
                                       {\ifthenelse{\isempty{#2}}{}{#2}}
                                       {\ifthenelse{\isempty{#2}}{#1}{#1  | #2}}}
\newcommand {\prmCat}   [2] {#1 \mathrel{\circ} #2}
\newcommand {\support}  [1] {\mathsf{supp}(#1)}
\newcommand {\cent}     [0] {\vdash}
\newcommand {\hsh}      [0] {\#}
\newcommand {\frsh}    [0] {\mathrel{\hsh}}
\newcommand {\freshS}  [3] {\mathsf{img^\hsh}(#1, \ #2, \ #3)}


\newcommand {\atms}     [1] {\mathsf{atms}(#1)}
\newcommand {\vars}     [1] {\mathsf{vars}(#1)}

\newcommand {\map}      [2] {[#1 \mapsto #2]}
\newcommand {\dom}      [1] {\mathsf{dom}(#1)}
\newcommand {\img}      [1] {\mathsf{img}(#1)}
\newcommand {\ds}       [1] {\mathsf{ds}(#1)}
\newcommand {\trans}    [2] {(#1 \ #2)}

\newcommand {\al}       [0] {{\approx_\alpha}}
\newcommand {\aleq}     [0] {\mathrel{\al}}
\newcommand {\aleqm}    [0] {\mathrel{\leftidx{^{?}}{\al}{}}}
\newcommand {\alequ}    [0] {\mathrel{\leftidx{^{?}}{\al}{^{?}}}}
\newcommand{\alphaEq}[3]{#1 \cent #2 \aleq #3}

\newcommand {\tSub}      [2] {\map{#1}{#2}}
\newcommand {\insVar}    [2] {\val{#1}{#2}}
\newcommand {\insAct}    [2] {#1 \, #2}
\newcommand {\joinTsubs} [2] {#1 \circ #2}

\newcommand {\vSub}     [2] {\map{#1}{#2}}
\newcommand {\subVar}   [2] {\val{#1}{#2}}
\newcommand {\subAct}   [2] {#1 \, #2}

\newcommand{\frshRuleName}[1]{\ensuremath{(\mathit{#1})^\hsh}}
\newcommand{\alphaRuleName}[1]{\ensuremath{(\mathit{#1})^\alpha}}

\newcommand{\concat}{+\kern-0.8em+}

\newcommand{\Trm}{\texttt{Term}}
\newcommand{\Frm}{\texttt{Form}}
\newcommand{\enc}{\texttt{enc}}
\newcommand{\dec}{\texttt{dec}}
\newcommand{\deH}{\ensuremath{\delta\enc_H}}
\newcommand{\dta}{\texttt{data}}
\newcommand{\LATy}{\ensuremath{\Lambda^{\rightarrow}\texttt{Type}}}
\newcommand{\LPTy}{\ensuremath{\Lambda^{\pi}\texttt{Type}}}
\newcommand{\LTrm}{\ensuremath{\Lambda\texttt{Term}}}
\newcommand{\LPTrm}{\ensuremath{\Lambda^{\pi}\texttt{Term}}}
\newcommand{\HOLT}{\ensuremath{\texttt{HOL}_\texttt{type}}}
\newcommand{\Exp}{\texttt{Exp}}
\newcommand{\SigmaEq}{\Sigma_\mathsf{eq}}
\newcommand{\SigmaNat}{\Sigma_\mathsf{nat}}



